% =====================================================================
% Section 3.2 (continued) — Two instantiations of the statistical prior
% and Section 3.3 — Personalized Tool Selection via Statistical Priors
% =====================================================================
% This file assumes the preamble already loads: amsmath, amssymb, algorithm,
% algpseudocode (or algorithm2e), and that the notation introduced in §3.1
% is in scope.
% =====================================================================


% ---------------------------------------------------------------------
% §3.2 continued — the two priors
% ---------------------------------------------------------------------

\paragraph{Frequency Prior.}
The simplest realisation of the local harness is a per-user empirical
success-rate table. For every triple $(u, d, t)$ comprising a user, a
domain, and a candidate tool, the harness maintains two scalars: the
number of attempts $N_{u,d,t}$ and the number of successes
$S_{u,d,t}\!=\!\sum_{\tau\le t}\mathbb{1}\{a_\tau\!=\!t\wedge r_\tau\!=\!1\}$.
Its preference distribution over the inferred domain's
candidate set $\mathcal{T}_{\hat d(q_t)}$ is the normalised success rate
\begin{equation}
\hat p^{\mathrm{freq}}_{u,d}(t) \;=\;
\frac{S_{u,d,t} / \max(N_{u,d,t},\,1)}
     {\sum_{t'\in\mathcal{T}_d} S_{u,d,t'} / \max(N_{u,d,t'},\,1)},
\label{eq:freq_prior}
\end{equation}
and its default action is the corresponding $\arg\max$. Because the
estimator is a pure sufficient statistic on the binary reward stream,
the update is $O(1)$ in time and memory, and the harness has no learnable
parameters beyond the counts themselves. The frequency prior makes no
use of the query text and performs no explicit exploration: untried
tools are visited only through an initial round-robin sweep over
$\mathcal{T}_d$. It therefore characterises the \emph{weakest possible}
local statistical primitive, and serves as a controlled point of
comparison for richer harnesses that account for context or
uncertainty.

\paragraph{Bandit Prior.}
A more principled realisation casts each $(u, d, t)$ as a separate
contextual-bandit arm. We instantiate this with
\textsc{LinUCB}~\citep{li2010linucb}: every arm maintains a $D$-dimensional
parameter $\boldsymbol{\theta}_{u,d,t}$ together with a regularised
Gram matrix $\mathbf{A}_{u,d,t}\in\mathbb{R}^{D\times D}$ and a moment
vector $\mathbf{b}_{u,d,t}\in\mathbb{R}^{D}$, updated online as
\begin{equation}
\mathbf{A} \leftarrow \mathbf{A} + \boldsymbol{\phi}(q,t,d)\boldsymbol{\phi}(q,t,d)^\top,
\quad
\mathbf{b} \leftarrow \mathbf{b} + r\,\boldsymbol{\phi}(q,t,d),
\quad
\boldsymbol{\theta} = \mathbf{A}^{-1}\mathbf{b}.
\label{eq:linucb_update}
\end{equation}
The context vector $\boldsymbol{\phi}(q,t,d)\in\mathbb{R}^{D}$ is
constructed by feature-hashing the query, the tool name, and the inferred
domain, producing a deterministic, vocabulary-free representation that
generalises across queries with related surface forms. At decision time
the harness computes an upper-confidence-bound score
\begin{equation}
s_{u,d,t}(q) \;=\;
\boldsymbol{\theta}_{u,d,t}^{\top}\boldsymbol{\phi}(q,t,d)
\;+\;
\alpha_{\mathrm{ucb}}\,
\sqrt{\boldsymbol{\phi}(q,t,d)^\top \mathbf{A}_{u,d,t}^{-1} \boldsymbol{\phi}(q,t,d)},
\label{eq:ucb_score}
\end{equation}
and exposes them either as a default action $\arg\max_t s_{u,d,t}(q)$ or
as a tempered-softmax preference distribution
$\hat p^{\mathrm{bandit}}_{u,d}(t)\propto\exp(\beta s_{u,d,t}(q))$.
The exploration coefficient $\alpha_{\mathrm{ucb}}$ controls the balance
between exploiting the current point estimate and probing under-sampled
arms. Compared with the frequency prior, the bandit prior is sensitive to
query context (through $\boldsymbol{\phi}$), and, more importantly,
maintains \emph{calibrated uncertainty}: the second term of
Equation~\eqref{eq:ucb_score} grows with epistemic uncertainty and shrinks
as data accumulate. Both properties become consequential when the user's
underlying preference distribution $\pi_u(d)$ is stochastic rather than
deterministic.

We emphasise that neither of these two priors is itself a contribution
of this work. Both are well-established estimators, and many alternative
local-statistical components (e.g.\ Thompson sampling, kernel bandits,
or a small neural recommender) are equally compatible with the
architecture. Our claim is that \emph{any} such local harness, when
combined with the override channel described next, is sufficient for
the personalised-tool-selection problem; the goal of comparing two
instantiations is to characterise how the strength of the local
statistical primitive interacts with the rest of the system, not to
identify a single optimal one.


% ---------------------------------------------------------------------
% §3.3 — overall method
% ---------------------------------------------------------------------

\subsection{Personalized Tool Selection via Statistical Priors}
\label{sec:method:overall}

We now describe the end-to-end decision procedure that the
\textsc{Local Harness} architecture induces, given a chosen statistical
prior $h\in\{\mathrm{freq},\mathrm{bandit}\}$. The procedure separates
each round of personalised tool selection into three sequential steps,
two of which are deterministic and local, and only one of which
optionally consults the remote language model.

\paragraph{Step~1: Shared domain classification.}
Upon receiving a query $q_t$, the system performs a single LLM call to
infer the target domain $\hat d(q_t)\in\mathcal{D}$. This call is
amortised across the entire agent stack rather than executed
per-component, so the per-round LLM cost does not scale with the
number of personalised hypotheses being evaluated. The resulting label
$\hat d(q_t)$ restricts the candidate set to
$\mathcal{T}_{\hat d(q_t)}\subset\mathcal{T}$ for all downstream
operations.

\paragraph{Step~2: Local statistical default.}
The local harness consumes the inferred domain together with the
user's accumulated state and returns a candidate action
\begin{equation}
\tilde a_t \;=\; \arg\max_{t\in\mathcal{T}_{\hat d(q_t)}}
\hat p^{\,h}_{u,\hat d(q_t)}(t),
\label{eq:harness_default}
\end{equation}
where $\hat p^{\,h}$ is either the success-rate distribution of
Equation~\eqref{eq:freq_prior} or the bandit-derived posterior implied
by Equation~\eqref{eq:ucb_score}. This step is fully local, deterministic
conditional on the user state, and incurs no LLM latency.

\paragraph{Step~3: Semantic override probe.}
The remote LLM is engaged only to check whether the user's query
contains an explicit lexical instruction that supersedes the user's
habitual preference. We issue a single binary probe (Appendix~G)
\begin{equation}
o_t,\,\tau_t \;=\; \mathrm{LLM}_{\text{override}}\!\left(q_t,\,
\mathcal{T}_{\hat d(q_t)}\right),
\label{eq:override_probe}
\end{equation}
where $o_t\in\{0,1\}$ indicates whether the query explicitly names a
specific tool and $\tau_t\in\mathcal{T}_{\hat d(q_t)}$ is the named
tool when $o_t=1$. The probe is deliberately narrow: it lists only the
tool \emph{names} (no descriptions, no statistical priors, no
chain-of-thought elicitation), reducing the language model's task to
lexical recognition rather than open-ended selection.

\paragraph{Action and update.}
The action executed at round $t$ is then
\begin{equation}
a_t \;=\;
\begin{cases}
\tau_t & \text{if } o_t = 1,\\[2pt]
\tilde a_t & \text{otherwise.}
\end{cases}
\label{eq:final_action}
\end{equation}
The environment returns the binary reward
$r_t=\mathbb{1}\{a_t=a^*_t\}$, after which the local harness alone is
updated with the tuple $(q_t,\hat d(q_t),a_t,r_t)$; the LLM is
stateless across rounds, by design. Compactly, given a choice of
prior $h\!\in\!\{\mathrm{freq},\mathrm{bandit}\}$ initialised with zero
state for every $(u,d,t)$, each round executes the pipeline
$q_t\!\to\!\hat d\!=\!\mathrm{LLM}_{\text{domain}}(q_t)
\!\to\!\tilde a\!=\!\arg\max_{t\in\mathcal{T}_{\hat d}}\hat p^{\,h}_{u,\hat d}(t)
\!\to\!(o,\tau)\!=\!\mathrm{LLM}_{\text{override}}(q_t,\mathcal{T}_{\hat d})
\!\to\! a_t\!=\!\tau\text{ if }o\!=\!1\text{ else }\tilde a
\!\to\! r_t\!=\!\mathbb{1}\{a_t\!=\!a^*_t\}
\!\to\!\text{Update}(h)$.

Several design properties of the procedure are worth making explicit.
First, the only path that involves the remote LLM in the regular
critical path of an interaction is the binary probe in
Equation~\eqref{eq:override_probe}, whose prompt is a fixed-length
template independent of the user's history. The number of remote-model
calls per round is therefore constant in both the size of the candidate
set and the amount of accumulated user state, sidestepping the
``lost-in-the-middle'' and quadratic-attention pathologies that
afflict memory-augmented baselines as the history grows. Second, the
local harness is the \emph{only} component whose state evolves across
rounds; the LLM is reset between calls. This separation of stateful
local learning from stateless remote reasoning is exactly what we
denote by ``decoupling'' and is the principal architectural claim of
this paper. Third, the choice of $h$ is entirely orthogonal to the
override mechanism: any local estimator that produces a probability
vector over $\mathcal{T}_d$ from the user's history can be substituted
into Step~2 without altering Steps~1 or~3. The two instantiations
introduced in \S\ref{sec:method:harness} are therefore not separate
methods, but two ways of populating the same architectural slot, and
serve in \S\ref{sec:exp} as controlled probes of how the strength of
the local statistical primitive affects overall performance.
